\documentclass[12pt]{article}
\usepackage[utf8]{inputenc}
\usepackage[T1]{fontenc}
\usepackage[margin=1in]{geometry}
\usepackage{setspace}
\usepackage{xcolor}
\usepackage[colorlinks=true, urlcolor=blue!60!black, linkcolor=blue!60!black]{hyperref}
\setlength{\parindent}{2em}
\setlength{\parskip}{1.25em}
\renewcommand{\baselinestretch}{1.0}
\begin{document}
\begin{center}
{\LARGE \textit{Jacobson v. Massachusetts}, 197 U.S. 11 (1905)\par}
\vspace{0.8em}
{\large Case Brief --- Four-Page Edition\par}
\end{center}
\vspace{1em}

\section{Memory Jogger}

During a 1902 Cambridge smallpox outbreak, the Court upheld Massachusetts's compulsory adult vaccination law and its \$5 fine as a valid exercise of the police power: Fourteenth Amendment liberty is ``liberty regulated by law,'' subject to reasonable restraints for the common good (\href{https://supreme.justia.com/cases/federal/us/197/11/}{Justia}).

\section{Detailed Case Facts}

A Massachusetts statute authorized a city's board of health, if ``in its opinion, it is necessary for the public health or safety,'' to require vaccination of all inhabitants, free of charge. An adult who refused forfeited \$5; the statute exempted only ``children who present a certificate, signed by a registered physician that they are unfit subjects for vaccination.''

In February 1902, citing the continuing spread of smallpox, the Cambridge Board of Health ordered vaccination of all inhabitants not vaccinated since 1897. Henning Jacobson, a Lutheran pastor over twenty-one, refused free vaccination. He said a childhood vaccination in Sweden had made him severely ill and that his son had suffered a similar reaction, but he offered no proof that, as an adult in 1902, he was medically unfit for vaccination.

\section{Procedural History}

Massachusetts prosecuted Jacobson in July 1902. On appeal to the Middlesex Superior Court, the trial judge excluded all fourteen of his offers of proof contesting vaccine safety and efficacy and refused his constitutional jury instructions; the jury convicted. The Supreme Judicial Court overruled his exceptions, and he was fined \$5. On writ of error, the U.S. Supreme Court heard argument on December 6, 1904, and decided the case on February 20, 1905.

\section{Judicial Votes}

Affirmed, \textbf{7--2}.
\begin{itemize}
\item \textbf{Majority:} Justice Harlan, joined by Chief Justice Fuller and Justices Brown, White, McKenna, Holmes, and Day.
\item \textbf{Concurrences:} None.
\item \textbf{Dissents:} Justices Brewer and Peckham, without written opinion.
\end{itemize}

\section{Holding}

Under its police power a State may enact reasonable public-health regulations, including authorizing local boards to require adult vaccination during an outbreak. Courts may invalidate such a law only if it ``has no real or substantial relation to those objects, or is, beyond all question, a plain, palpable invasion of rights secured by the fundamental law.'' Exempting medically unfit children but not adults does not deny equal protection. The Court decided ``only that the statute covers the present case''; its warnings about adults whose health vaccination would ``seriously impair'' are dicta.

\section{Analysis of Opinions}

\textbf{Majority (Harlan).} The Preamble confers no substantive rights, and the Constitution's ``spirit'' ``is to be collected chiefly from its words.'' State authority rests on the reserved police power, which includes ``quarantine laws and `health laws of every description''' and on the social compact. Liberty does not mean being wholly free from restraint; quarantine of healthy travelers and military conscription show that the common good can override individual choice. Because medical consensus supported vaccination, choosing among competing medical theories belonged to the legislature, not a jury. Under \textit{Mugler v. Kansas}, courts intervene only where a health law bears no real or substantial relation to its object or plainly invades fundamental rights. ``Obvious reasons'' distinguished children from adults. But police power may not be exercised in an ``arbitrary and oppressive'' manner, and forcing vaccination on an adult for whom it would be life-threatening would be ``cruel and inhuman in the last degree.''

\textbf{Dissents (Brewer and Peckham).} They wrote nothing, but two months later Peckham, joined by Brewer, wrote \textit{Lochner v. New York}, striking down a bakery-hours law, while Harlan dissented there by invoking \textit{Jacobson}'s deferential police-power framework.

\section{Examples: Future Applications}

\begin{enumerate}
\item \textbf{Same-Side (Upheld): School Immunization Mandate.} Barring unvaccinated children from school unless medically exempt is a reasonable disease-control measure; the Court applied \textit{Jacobson} to uphold one in \textit{Zucht v. King} (1922).
\item \textbf{Same-Side (Upheld): Quarantine of Exposed Travelers.} Isolating arriving passengers for fourteen days during a hemorrhagic-fever outbreak tracks \textit{Jacobson}'s approval of maritime quarantine.
\item \textbf{Opposite-Side (Invalid): Medically Contraindicated Adult.} Fining and threatening to jail an adult with documented severe anaphylaxis falls within the closing reservation: vaccination that would ``seriously impair his health or probably cause his death'' is arbitrary and oppressive as applied.
\item \textbf{Opposite-Side (Invalid): Singling Out Worship.} Capping religious services at ten people while sparing comparable secular venues fails; \textit{Roman Catholic Diocese of Brooklyn v. Cuomo} (2020) held \textit{Jacobson} does not suspend heightened scrutiny for First Amendment rights.
\item \textbf{Fence-Sitter: Indefinite Booster Fines.} A permanent \$500-per-month fine for declining an annual booster against a low-fatality endemic virus: supporters invoke legislative deference; challengers stress that \textit{Jacobson} involved an acute epidemic, a \$5 fine, and a vaccine that stopped transmission.
\end{enumerate}

\section{Critique}

\textbf{Originalist and libertarian critics.} Justice Gorsuch in \textit{Cuomo} and scholars such as Josh Blackman argue that \textit{Jacobson} was a modest, pre-incorporation decision applying early rational-basis review to a \$5 fine, not a ``super-precedent'' suspending the Bill of Rights in emergencies. Critics across the spectrum stress its darkest legacy: Holmes cited it as the sole precedent for compulsory sterilization in \textit{Buck v. Bell}.

\textbf{Public-health and progressive views.} Lawrence Gostin reads \textit{Jacobson} as a counterweight to \textit{Lochner} built on necessity, reasonable means, proportionality, and harm avoidance. Wendy Mariner and colleagues caution that it predates incorporation and strict scrutiny, and should not license arbitrary intrusions on bodily integrity.

\section{Key Quotations}

\begin{quote}
``But the liberty secured by the Constitution of the United States to every person within its jurisdiction does not import an absolute right in each person to be, at all times and in all circumstances, wholly freed from restraint. There are manifold restraints to which every person is necessarily subject for the common good.''
\end{quote}
\noindent --- Justice Harlan, for the Court, 197 U.S. at 26.

\begin{quote}
``Upon the principle of self-defense, of paramount necessity, a community has the right to protect itself against an epidemic of disease which threatens the safety of its members.''
\end{quote}
\noindent --- Justice Harlan, for the Court, 197 U.S. at 27.

\begin{quote}
``[T]he police power of a State, whether exercised by the legislature or by a local body acting under its authority, may be exerted in such circumstances or by regulations so arbitrary and oppressive in particular cases as to justify the interference of the courts to prevent wrong and oppression.''
\end{quote}
\noindent --- Justice Harlan, for the Court, 197 U.S. at 38.

\end{document}
